Une caisse de masse $m=\qty{20.0}{\kg}$ est tirée sur un sol horizontal supposé
parfaitement lisse (absence de frottement). Le câble de traction fait un angle
$\alpha=\ang{60}$ avec l'horizontale et la force de traction a pour valeur
$F=\qty{10}{\N}$.
\begin{enumerate}
    \item Représenter les forces s'exerçant sur la caisse sur un schéma. Les
          nommer et préciser leurs caractéristiques.
    \item Calculer le travail de chacune des forces lorsque la caisse se déplace
          de $\qty{5.0}{\m}$ sur le sol.
    \item Reprendre les questions précédentes en supposant que le sol est
          rugueux (existence de frottements), la valeur de la force de
          frottement étant $f=\qty{0.80}{\N}$.
\end{enumerate}
